% DO NOT EDIT - automatically generated from metadata.yaml

\def \codeURL{[public repository URL]}
\def \codeDOI{}
\def \codeSWH{}
\def \dataURL{https://github.com/GATECH-EIC/HW-NAS-Bench}
\def \dataDOI{}
\def \editorNAME{}
\def \editorORCID{}
\def \reviewerINAME{}
\def \reviewerIORCID{}
\def \reviewerIINAME{}
\def \reviewerIIORCID{}
\def \dateRECEIVED{}
\def \dateACCEPTED{}
\def \datePUBLISHED{}
\def \articleTITLE{[¬Re] LLM-NAS: LLM-Driven Hardware-Aware Neural Architecture Search}
\def \articleTYPE{Replication}
\def \articleDOMAIN{Machine Learning}
\def \articleBIBLIOGRAPHY{bibliography.bib}
\def \articleYEAR{2026}
\def \reviewURL{}
\def \articleABSTRACT{LLM-NAS uses a large language model (LLM) as the variation operator of a hardware-aware neural architecture search and reports near-optimal Pareto fronts on HW-NAS-Bench from 120 calls to the language model per search. No reference implementation was available, so we reimplemented the method from the paper's text, algorithm and prompts in a configuration-driven framework that logs every call for exact replay. We first validated the benchmark data: eleven of the paper's twelve reported architectures were recovered exactly, one reported point does not exist in the benchmark, and 4,699 architectures carry an invalid Edge TPU latency. We then compared searches driven by GPT-4.1, Claude Opus 5.5, Claude Haiku 4.5, DeepSeek-Coder-V2 and Qwen2.5-7B with non-LLM controls that share the evaluator and the niche structure, using hypervolume (HV) and inverted generational distance. Under the default convention GPT-4.1 reached HV 0.964 (paper: 0.997) and the controls 0.989–0.991. The ranking depends on the normalisation: one degenerate architecture on the true front compresses all sensible architectures into 9\% of the accuracy axis, and excluding architectures below 80\% accuracy places Claude Opus 5.5 first (0.961) and GPT-4.1 below the controls (0.752). Across six devices, none of 24 LLM runs exceeded random-operator evolution and one exceeded random sampling. On this small tabulated space with a cheap evaluator, the method's cost advantage comes from the training-free evaluator rather than from the language model; whether the LLM operator pays off with expensive evaluations or in larger spaces remains open.}
\def \replicationCITE{Zhu, H., Zhang, G. L. and Huang, S. LLM-NAS: LLM-driven Hardware-Aware Neural Architecture Search. arXiv:2510.01472, 2025.}
\def \replicationBIB{llmnas}
\def \replicationURL{https://arxiv.org/abs/2510.01472}
\def \replicationDOI{10.48550/arXiv.2510.01472}
\def \contactNAME{Siyabonga Zimozoxolo Mjali}
\def \contactEMAIL{u12167283@up.ac.za}
\def \articleKEYWORDS{neural architecture search, hardware-aware nas, large language models, reproducibility, evolutionary computation, multi-objective optimisation, python}
\def \journalNAME{ReScience C}
\def \journalVOLUME{}
\def \journalISSUE{}
\def \articleNUMBER{}
\def \articleDOI{}
\def \authorsFULL{Siyabonga Zimozoxolo Mjali and Thambo Nyathi}
\def \authorsABBRV{S.Z. Mjali and T. Nyathi}
\def \authorsSHORT{Mjali and Nyathi}
\title{\articleTITLE}
\date{}
\author[1,\orcid{0009-0008-9207-5371}]{Siyabonga Zimozoxolo Mjali}
\author[1,\orcid{0000-0003-4676-6063}]{Thambo Nyathi}
\affil[1]{Department of Computer Science, University of Pretoria, Corner of Lynnwood Road and Roper Street, Hatfield, Pretoria 0002, South Africa}
